% XeLaTeX and Traditional Chinese typesetting.
% These are working layout settings, ready to adapt to a university template.
\usepackage{iftex}
\ifXeTeX\else
  \PackageError{thesis-abstract}{Select XeLaTeX as the compiler}
    {In Overleaf, set Compiler to XeLaTeX and Main document to main.tex.}
\fi

\usepackage[left=30mm,right=25mm,top=25mm,bottom=25mm]{geometry}
\usepackage{fontspec}
\usepackage{setspace}
\usepackage[unicode,hidelinks]{hyperref}

% Fonts available on Overleaf; no Windows-only fonts are required.
\setmainfont{texgyretermes}[
  Extension = .otf,
  UprightFont = *-regular,
  BoldFont = *-bold,
  ItalicFont = *-italic,
  BoldItalicFont = *-bolditalic
]
\setsansfont{texgyreheros}[
  Extension = .otf,
  UprightFont = *-regular,
  BoldFont = *-bold,
  ItalicFont = *-italic,
  BoldItalicFont = *-bolditalic
]
\setmonofont{lmmono10-regular.otf}[
  BoldFont = lmmonolt10-bold.otf,
  ItalicFont = lmmono10-italic.otf,
  BoldItalicFont = lmmonolt10-boldoblique.otf
]
\setCJKmainfont{Noto Serif CJK TC}
\setCJKsansfont{Noto Sans CJK TC}
\setCJKmonofont{Noto Sans CJK TC}

\ctexset{
  chapter = {
    format = \centering\bfseries\Large,
    beforeskip = 0pt,
    afterskip = 24pt
  }
}
\setstretch{1.3}
\setlength{\parindent}{2em}
\setlength{\parskip}{0.3\baselineskip}
\setlength{\emergencystretch}{2em}

\hypersetup{
  pdftitle = {動態環境下移動機械臂預測導航與受約束操作},
  pdfsubject = {中英文摘要與關鍵字}
}

% Hanging layout for wrapped keyword lines, shared by both abstracts.
\newlength{\thesiskeywordlabelwidth}
\newcommand{\thesiskeywords}[2]{%
  \par\vspace{0.8\baselineskip}%
  \begingroup
  \settowidth{\thesiskeywordlabelwidth}{\textbf{#1}}%
  \hangindent=\thesiskeywordlabelwidth
  \hangafter=1
  \noindent\textbf{#1}#2\par
  \endgroup
}
